
This work provides evidence that how error information is organized affects LLM self-repair success, independent of what information is present. The key finding is that structured error formatting outperforms scrambled formatting with identical content ($\Delta$ = +14.4\%, p < 0.001).

The fix specificity experiment (simulation mode) shows an inverted-U pattern consistent with scaffolding theory, with intermediate-level hints achieving highest success rates. The scale interaction test finds a directional pattern but does not reach statistical significance.

These findings suggest that error format organization is a relevant design variable for self-repair systems. A simple preprocessing step---reformatting compiler errors with section labels---may improve self-repair without model changes.

\subsubsection{Future Work}

Immediate extensions include: (1) real-model validation for H-M3 by resolving flash\_attn compatibility; (2) API-based validation for H-M1 reconstruction test; (3) increased sample sizes for scale interaction analysis.

Longer-term directions include: cross-language validation (Java, JavaScript, C++), evaluation on production codebases (SWE-bench), and investigation of whether optimal formats can be learned automatically.
